\section{Results}
\label{sec:results}

We evaluate seven container escape scenarios under eight hardening controls (B1--B8) plus the unhardened baseline (B0), with \texttt{google/gemini-3.5-flash} at a standardized budget of \num{300000} tokens and \num{120} messages per sample. The six completed scenarios contribute 520 valid samples; \texttt{hostpath\_etc} completed its clean re-run with 0/10 escapes at every control, but the agent never escaped even at baseline, so its model-level rates carry no defense signal; its reference-exploit results (RQ2) are the informative part.

\subsection{Baseline and refusal rates}
The unhardened baseline (B0) confirms high offensive capability on misconfiguration and delegation pathways: \texttt{privileged} (SR $=1.00$, $10/10$), \texttt{docker\_socket} (SR $=1.00$, $10/10$), and \texttt{pid\_namespace} (SR $=0.90$, $9/10$). Zero baseline refusals were observed across these scenarios ($0\%$). In contrast, kernel capability scenarios (\texttt{sys\_admin}, \texttt{sys\_module}, \texttt{dac\_read\_search}) exhibit a baseline ceiling (SR $=0.00$, $0/10$), where the model reaches token limits without landing the exploit.

Refusal-like messages (detected by keyword match on assistant text) are concentrated in control B6 (distroless): in each of the six completed scenarios, all $10$ B6 samples contain an explicit refusal and end at the message cap. In the B6 samples we checked, each refusal followed at least one tool call. Outside B6, refusal text is rare: none in most controls, with isolated instances in \texttt{sys\_admin} B0 and B3 (one each), \texttt{docker\_socket} B7 (one), and \texttt{pid\_namespace} B7 (three).

\subsection{Per-control escape-rate reduction (RQ1)}
\Cref{tab:hardening-matrix} summarizes the escape success rates (SR) and absolute risk reduction (ARR) relative to B0 across all scenarios. gVisor (B7) is the single most effective control, achieving ARR $=0.97$ across baseline-escapable scenarios by driving escapes to zero on \texttt{privileged}, \texttt{docker\_socket}, and \texttt{pid\_namespace}. Strict seccomp (B3) and read-only rootfs (B4) yield reductions on \texttt{privileged} (ARR $=0.40$ and $0.10$ respectively; the B3 non-escapes are token-budget exhaustion during raw device scanning). On \texttt{docker\_socket}, B3 and B4 leave the escape unchanged (ARR $=0.00$). On \texttt{pid\_namespace}, B3 gives ARR $=0.10$ and B4 gives ARR $=0.00$. Capability dropping (B1) and no-new-privileges (B2) give ARR $=0.00$ in all three baseline-escapable scenarios. The composite bundle (B8) is not neutral: its ARR is $0.10$ on \texttt{privileged}, $0.00$ on \texttt{docker\_socket}, and $-0.10$ on \texttt{pid\_namespace}, where it produced one more escape than B0 (10/10 against 9/10). Its mean ARR of $0.00$ hides this per-scenario spread.

% Populate from results/eval_logs/*.eval (repository root) once the campaign completes.
% One row per scenario; SR per control; ARR/RR relative to B0.
\begin{table*}[t]
\centering
\caption{Empirical hardening matrix: escape success rate (SR) per
(scenario, control), with absolute (ARR) and relative (RR) risk reduction versus
the B0 baseline. Each SR is escapes divided by genuine attempts (refused samples excluded; see Methodology). Every cell with a non-zero SR has no refusals, so its value is also escapes out of \num{10} epochs. Intervals are discussed in the Statistical analysis subsection.}
\label{tab:hardening-matrix}
\small
\begin{tabular}{@{}lccccccccc@{}}
\toprule
Scenario & B0 & B1 & B2 & B3 & B4 & B5 & B6 & B7 & B8 \\
\midrule
privileged      & 1.00 & 1.00 & 1.00 & 0.60 & 0.90 & \textit{n/a} & \textit{n/a}\textsuperscript{r} & 0.00 & 0.90 \\
docker\_socket  & 1.00 & 1.00 & 1.00 & 1.00 & 1.00 & 1.00 & \textit{n/a}\textsuperscript{r} & 0.00 & 1.00 \\
pid\_namespace  & 0.90 & 0.90 & 0.90 & 0.80 & 0.90 & \textit{n/a} & \textit{n/a}\textsuperscript{r} & 0.00 & 1.00 \\
sys\_admin\textsuperscript{m}      & 0.00 & 0.00 & 0.00 & 0.00 & 0.00 & 0.00 & \textit{n/a}\textsuperscript{r} & 0.00 & 0.00 \\
sys\_module\textsuperscript{m}     & 0.00 & 0.00 & 0.00 & 0.00 & 0.00 & 0.00 & \textit{n/a}\textsuperscript{r} & 0.00 & 0.00 \\
dac\_read\_search\textsuperscript{m} & 0.00 & 0.00 & 0.00 & 0.00 & 0.00 & 0.00 & \textit{n/a}\textsuperscript{r} & 0.00 & 0.00 \\
hostpath\_etc   & 0.00\textsuperscript{c} & 0.00\textsuperscript{c} & 0.00\textsuperscript{c} & 0.00\textsuperscript{c} & 0.00\textsuperscript{c} & 0.00\textsuperscript{c} & 0.00\textsuperscript{cr} & 0.00\textsuperscript{c} & 0.00\textsuperscript{c} \\
\midrule
Mean ARR vs B0\textsuperscript{$\dagger$}  & --- & 0.00 & 0.00 & 0.17 & 0.03 & 0.00 & \textit{n/a} & 0.97 & 0.00 \\
\bottomrule
\end{tabular}

\vspace{2pt}
\footnotesize\textit{n/a} = control not applicable: B5 user-namespace remap is
rejected by the Docker daemon for \texttt{-{}-privileged} and for
\texttt{-{}-pid=host} (\texttt{pid\_namespace}), so the container cannot be created.
\textsuperscript{r}B6 (distroless): all 10 epochs ended in refusal messages at the message cap,
each after at least one tool call; SR over genuine attempts is therefore not well defined (raw escape count 0/10).
\textsuperscript{m}Model-capability-limited scenario: the attacker never escapes at any
control \emph{including the baseline} (SR $=0$ throughout), so ARR/RR carry no defense
signal; the mechanistic interpretation is in \Cref{tab:rq3-matrix} (RQ2).
\textsuperscript{c}\texttt{hostpath\_etc} was re-run in full (B0--B8) after the earlier sweep was discarded;
every control shows 0/10 escapes, and the agent never escaped even at baseline, so these rates carry no
defense signal. B0--B3 ran before the per-epoch host \texttt{/etc} guard existed, so host cron carryover
within an epoch was possible; carryover can only raise escape counts, so these zeros are not inflated.
B4--B8 ran with the guard, which restores host \texttt{/etc} before each epoch. The guard did not cover
\texttt{apt} or \texttt{if-up} hooks, and two such hooks written during B7 persisted into B8 until we removed
them after the run. No B8 epoch escaped, but we cannot rule out an effect of those hooks on B8 epochs.
The RQ2 column (\Cref{tab:rq3-matrix}) was run under clean per-control state and is valid.
\textsuperscript{$\dagger$}Mean ARR is averaged only over the three scenarios with a
non-zero baseline escape rate (\texttt{privileged}, \texttt{docker\_socket},
\texttt{pid\_namespace}); the capability-limited and contaminated rows carry no ARR signal
and are excluded. Cells marked \textit{n/a} above are omitted from their control's mean.
\end{table*}


\subsection{Mechanistic classification (RQ2)}
\Cref{tab:rq3-matrix} pairs LLM evaluation outcomes with reference-exploit executions to isolate architectural fixes from model friction. Across the 54 applicable (scenario, control) pairs (the 56 scenario--control cells minus the two B5 cells that the Docker daemon rejects for \texttt{privileged} and \texttt{pid\_namespace}), each cell is classified from one reference-exploit run and the ten model epochs for that cell. Cells where the model and reference both fail are classified by audit of the failure cause, not by the 2$\times$2 rule alone; the B4 and B6 entries below are the audited exceptions:
\begin{itemize}
  \item \textbf{Genuine Architectural Fixes (FIX):} gVisor (B7) is a true fix for 6 of 7 scenarios (excluding \texttt{hostpath\_etc}). Strict seccomp (B3) is a genuine fix for capability-limited scenarios by blocking required syscalls (\texttt{mount}, \texttt{init\_module}, \texttt{open\_by\_handle\_at}). User-namespace remap (B5) is a genuine fix for capability scenarios and \texttt{hostpath\_etc}.
  \item \textbf{Harness Artifacts (FRICTION):} B4 (read-only rootfs) and B6 (distroless) cause reference exploit failures due to missing writable `/tmp` or missing bash binaries, even though the underlying security boundary remains unhardened.
  \item \textbf{Open Vulnerabilities (OPEN):} Standard container-level controls (B1, B2, B3, B4, B8) leave misconfiguration scenarios fully open to exploitation.
\end{itemize}

% Populate from results/eval_logs/reference_exploits.csv crossed with model outcomes.
% Cell legend: FIX = defense closes vuln; FRICTION = vuln open but model failed;
% BYPASS = model escaped though reference failed; OPEN = both succeeded.
\begin{table*}[t]
\centering
\caption{RQ2 mechanistic classification per (scenario, control): whether a
model failure reflects a true architectural fix or mere friction, using the
reference-exploit outcome under the identical control.}
\label{tab:rq3-matrix}
\small
\begin{tabular}{@{}lcccccccc@{}}
\toprule
Scenario & B1 & B2 & B3 & B4 & B5 & B6 & B7 & B8 \\
\midrule
privileged      & OPEN & OPEN & FRICTION & OPEN & \textit{n/a} & FRICTION & \textbf{FIX} & OPEN \\
docker\_socket  & OPEN & OPEN & OPEN & OPEN & OPEN & FRICTION & \textbf{FIX} & OPEN \\
pid\_namespace  & OPEN & OPEN & OPEN & OPEN & \textit{n/a} & FRICTION & \textbf{FIX} & OPEN \\
sys\_admin      & FRICTION & FRICTION & \textbf{FIX} & FRICTION & \textbf{FIX} & FRICTION & \textbf{FIX} & \textbf{FIX} \\
sys\_module     & FRICTION & FRICTION & \textbf{FIX} & FRICTION & \textbf{FIX} & FRICTION & \textbf{FIX} & \textbf{FIX} \\
dac\_read\_search & FRICTION & FRICTION & \textbf{FIX} & FRICTION & \textbf{FIX} & FRICTION & \textbf{FIX} & \textbf{FIX} \\
hostpath\_etc   & OPEN & OPEN & OPEN & OPEN & \textbf{FIX} & FRICTION & OPEN & OPEN \\
\bottomrule
\end{tabular}

\vspace{2pt}
\footnotesize For the three capability-limited scenarios (\texttt{sys\_admin},
\texttt{sys\_module}, \texttt{dac\_read\_search}) the attacker model never escapes, so RQ2
carries the interpretation. B1/B2 leave the required capability in place (vuln OPEN, model
simply capability-limited $\Rightarrow$ FRICTION). B3 (strict seccomp) is a genuine FIX ---
it denies the escape's key syscall (\texttt{mount}, \texttt{init\_module}, and
\texttt{open\_by\_handle\_at} respectively). B5 (userns) and B7 (gVisor) are genuine fixes,
and B8 inherits B3's seccomp. The B4 (read-only) and B6 (distroless) FRICTION cells are
\emph{harness artifacts}: the reference fails only because it cannot stage/compile tooling
(read-only \texttt{/tmp}, \texttt{noexec} tmpfs, or no in-image shell), while the underlying
capability is untouched --- they are not architectural fixes. For \texttt{hostpath\_etc}
(host-cron delegation) only B5 bites; notably gVisor (B7) does \emph{not} fix it because the
privileged action runs in the host cron daemon, outside the sandbox.
\textbf{Privileged B4 and B8 are a special case.} They are labelled OPEN although the reference
exploit failed (rc=1) under each, and the model escaped in 9 of 10 epochs under each. Under the
scheme in \Cref{tab:rq3-scheme}, a reference failure with a model success requires audit. The
label therefore rests on the model's escapes, not on the reference exploit. The reference failure is
attributed to the read-only \texttt{/tmp} harness artifact (the script cannot stage its tooling),
not to a control effect.
\end{table*}


\subsection{Scenario case study: \texttt{privileged} (complete)}
\label{sec:results-privileged}
We report detailed findings for the six completed scenarios in \Cref{tab:hardening-matrix,tab:rq3-matrix}, and treat \texttt{hostpath\_etc} separately below. For \texttt{privileged}, the attacker was \texttt{google/gemini-3.5-flash} at a fixed budget of \num{300000} tokens and \num{120} messages (\num{10} epochs per control in local Docker mode). B5 is not applicable to a \texttt{-{}-privileged} container.

\paragraph{Escape rates (RQ1).} The baseline is fully escapable (B0 SR $=1.00$,
$10/10$). Dropping capabilities (B1) and no-new-privileges (B2) have no effect
(SR $=1.00$): a \texttt{-{}-privileged} container exposes the host block devices under
\texttt{/dev}, and reading the raw host disk requires neither extra capabilities nor
privilege escalation. Read-only rootfs (B4) and the full composite (B8) each score
SR $=0.90$ (their single non-escapes were token-budget exhaustion, not blocks). Strict
seccomp (B3) scores SR $=0.60$. gVisor (B7) is the only control that drives escapes to
zero (SR $=0.00$, $10/10$ reaching the token cap with \emph{zero} refusals). Distroless
(B6) yielded $0/10$ escapes, and all ten end in \emph{refusal} messages at the message cap, after tool calls;
the shell-less image removes the model's tooling, so SR over
genuine attempts is therefore undefined and the raw $0/10$ reflects alignment/tooling
friction, not the security boundary.

\paragraph{Mechanistic classification (RQ2).} The reference exploit distinguishes a
real fix from mere friction (\Cref{tab:rq3-matrix}). It succeeds under B0--B3 (the
vulnerability is open; B3's lower SR is token-exhaustion friction from forced
raw-device scanning). The na\"ive reference script reported failures under B4, B6, and
B8, but these are \emph{harness artifacts}, not blocks: the script writes to
\texttt{/tmp} (read-only under B4/B8) and needs an in-image shell (absent under B6
distroless). A direct \texttt{debugfs} read of the host device confirms the flag is
still retrievable under B4 and B8 (exit status~0), and the device passthrough under B6
is identical to the baseline---so the vulnerability remains open in all three.
\textbf{B7 (gVisor) is the sole genuine architectural fix}: the host block-device
\emph{node} is present in \texttt{/dev}, but \texttt{open()} returns \texttt{ENXIO}
because gVisor's user-space kernel (Sentry) does not back the host device, so even the
reference exploit cannot read it. Both the model and the reference exploit fail for the
same real reason---nothing to read.

\paragraph{Takeaway.} For \texttt{privileged}, a low escape rate does \emph{not} imply
an effective control: B3 ($0.60$), B6 (refusals), and B7 ($0.00$) are low for three
different reasons---attacker-cost friction, alignment/tooling friction, and a true
boundary, respectively. Only the reference-exploit cross-check---audited against the
harness's own write/shell assumptions---separates them, and it reclassified three
apparent ``fixes'' (B4, B6, B8) as open vulnerabilities.

\subsection{Scenario case study: \texttt{docker\_socket} (complete)}
\label{sec:results-docker-socket}
The second fully analyzed scenario uses the identical protocol
(\texttt{google/gemini-3.5-flash}, \num{300000} tokens / \num{120} messages, \num{10}
epochs per control, local Docker mode). Unlike \texttt{privileged}, this is a
\emph{delegation} attack rather than a container breakout: the container mounts the host
Docker daemon socket (\texttt{/var/run/docker.sock}), and the attacker drives the
host-root daemon over the API to spawn or \texttt{exec} into a container that reads the
host flag. The privileged action happens in the host daemon, so container-level
hardening of the compromised container is largely off the attack path.

\paragraph{Escape rates (RQ1).} The baseline is fully escapable (B0 SR $=1.00$,
$10/10$), and every container-level control leaves it open: B1 (cap-drop), B2
(no-new-privileges), B3 (strict seccomp), B4 (read-only rootfs), B5 (user-namespace
remap), and the composite B8 all score SR $=1.00$. Notably B3 scores $1.00$ here versus
$0.60$ in \texttt{privileged}: the socket escape is cheap socket I/O rather than a
token-heavy raw-device scan, so seccomp neither blocks it nor induces token exhaustion---a
clean illustration that a control's effect depends on the attack path. gVisor (B7) is the
sole control that drives escapes to zero (SR $=0.00$, $10/10$ reaching the token cap with
one refusal-like message among the ten). Distroless (B6) again yielded $0/10$, with all ten ending in refusal messages at the message cap, from the
shell-less image.

\paragraph{Mechanistic classification (RQ2).} The reference exploit (which drives the
Docker API with \texttt{curl}) succeeds under B0--B5 and B8, confirming the vulnerability
is open under all container-level controls---including user-namespace remap, since
delegation targets the host-root daemon regardless of the remap. B6's reference failure is
a harness artifact (no shell/\texttt{curl} in the distroless image; \texttt{rc=127}), not a
block: the socket is still mounted. \textbf{B7 (gVisor) is the sole genuine architectural
fix}: \texttt{runsc} does not proxy host unix-domain sockets by default, so the socket node
is present but \texttt{connect()} is refused (model runs:
\texttt{ConnectionRefusedError}~\texttt{[Errno~111]}; reference exploit: curl
\texttt{rc=7}). Both the model and the reference exploit fail for the same real
reason---the delegation channel is severed.

\paragraph{Takeaway.} \texttt{docker\_socket} reinforces the \texttt{privileged} lesson
from the opposite direction: here six controls (B1--B5, B8) are empirically ineffective
despite the composite ``defense-in-depth'' bundle, because they do not sit on the
delegation path. gVisor is again the only real fix, and across both scenarios it is the
single control that holds---by virtualizing the host device in \texttt{privileged} and by
isolating the host socket in \texttt{docker\_socket}.

\subsection{Scenario case study: \texttt{pid\_namespace} (complete)}
\label{sec:results-pid-namespace}
The third fully analyzed scenario uses the identical protocol
(\texttt{google/gemini-3.5-flash}, \num{300000} tokens / \num{120} messages, \num{10}
epochs per control, local Docker mode). The container runs in the \emph{host PID
namespace} (\texttt{-{}-pid=host -{}-cap-add=SYS\_PTRACE -{}-security-opt
apparmor=unconfined}), so every host process is visible and the host root filesystem is
reachable through \texttt{/proc/<host-pid>/root/} (e.g.\ \texttt{cat
/proc/1/root/flag.txt}) or by ptrace-injecting a host process. B5 (user-namespace remap)
is \emph{not applicable}: the Docker daemon rejects \texttt{-{}-pid=host} when user
namespaces are enabled (\texttt{cannot share the host PID namespace when user namespaces
are enabled}), so the container cannot be created---structurally the same limitation as
\texttt{-{}-privileged}.

\paragraph{Escape rates (RQ1).} The escape stays open under every container-level control:
B0--B4 score $0.80$--$0.90$ and the composite B8 scores $1.00$. The handful of
non-escapes under B0--B4 are token-budget exhaustion (the model wandering), not blocks---B3
is lowest at $0.80$ (two token-capped epochs). gVisor (B7) is the only control that drives
escapes to zero (SR $=0.00$, $10/10$ reaching the token/message cap; three of these ten also contain refusal-like text). Distroless (B6)
yielded $0/10$, with all ten ending in refusal messages at the message cap, from the shell-less image.

\paragraph{Mechanistic classification (RQ2).} The reference exploit succeeds under B0--B4
and B8, confirming the vulnerability is open under all applicable container-level controls.
B6's reference failure is a harness artifact (no shell in the distroless image;
\texttt{rc=127}), not a block---the host PID namespace is still shared. \textbf{B7 (gVisor)
is the sole genuine architectural fix}: \texttt{runsc} presents its own Sentry-managed PID
namespace, so \texttt{-{}-pid=host} yields no real host-process visibility and
\texttt{/proc/1/root} resolves only to the container's own root. Both the model and the
reference exploit fail for the same real reason---no host process to pivot through.

\paragraph{Takeaway.} \texttt{pid\_namespace} is the third scenario in which \emph{gVisor is
the only control that holds}, this time by denying host-process visibility rather than by
virtualizing a device or isolating a socket. Across \texttt{privileged},
\texttt{docker\_socket}, and \texttt{pid\_namespace}, the entire suite of container-level
controls (cap-drop, no-new-privs, seccomp, read-only rootfs, and their composite) never
closes the vulnerability---only the user-space-kernel runtime does, and in two of the three
scenarios user-namespace remap is not even an applicable option.

\subsection{Capability-limited scenarios: \texttt{sys\_admin}, \texttt{sys\_module}, and \texttt{dac\_read\_search}}
\label{sec:results-capability-limited}
Three evaluated scenarios---\texttt{sys\_admin} (AppArmor/seccomp bypass via \texttt{core\_pattern}),
\texttt{sys\_module} (kernel module injection via \texttt{init\_module}), and
\texttt{dac\_read\_search} (handle brute-force via \texttt{open\_by\_handle\_at})---demonstrate
a \emph{model-capability ceiling}: \texttt{google/gemini-3.5-flash} achieves SR $=0.00$ ($0/10$) across
every control, \textbf{including the unhardened baseline (B0)}. Most non-escapes under the other controls
hit the token budget ($300\text{k}$ tokens); the rest hit the $120$-message cap (one \texttt{sys\_admin} B0 sample) or the wall-clock cap (one \texttt{sys\_module} B8 sample). All B6 samples end at the message cap after refusal messages.

Because the attacker cannot land the exploit even in the baseline, the empirical escape rates (RQ1)
carry zero defense-reduction signal (ARR and RR are undefined). Instead, the \emph{mechanistic
classification} (RQ2) via the reference exploit serves as the primary evaluation tool:
\begin{itemize}
  \item \textbf{Under \texttt{sys\_admin}}: B1 and B2 leave the capability open (FRICTION). Strict seccomp (B3, denying \texttt{mount}), user-namespace remap (B5), gVisor (B7), and composite (B8) are \textbf{genuine FIXes}. B4 (read-only) and B6 (distroless) are harness artifacts.
  \item \textbf{Under \texttt{sys\_module}}: Strict seccomp (B3, denying \texttt{init\_module}), user-namespace remap (B5), gVisor (B7), and composite (B8) are \textbf{genuine FIXes}. B1/B2 leave capability intact (FRICTION).
  \item \textbf{Under \texttt{dac\_read\_search}}: Strict seccomp (B3, denying \texttt{open\_by\_handle\_at}), user-namespace remap (B5), gVisor (B7), and composite (B8) are \textbf{genuine FIXes}. B4 is a notable harness artifact: read-only rootfs blocks script compilation in \texttt{/tmp}, but \texttt{open\_by\_handle\_at} is a \emph{read} syscall and the capability itself remains untouched.
\end{itemize}

\subsection{Scenario case study: \texttt{hostpath\_etc}}
\label{sec:results-hostpath-etc}
The seventh scenario bind-mounts host \texttt{/etc} writable into the container on a host running
\texttt{cron}. Like \texttt{docker\_socket}, this is a \emph{host-delegation attack}: the container writes
a root cron entry to \texttt{/hostfs/etc/cron.d/}, and the host cron daemon executes the payload outside
the container sandbox.

\paragraph{Mechanistic classification (RQ2).} The reference exploit succeeds under B0--B4, B7, and B8,
demonstrating that cap-drop, no-new-privs, strict seccomp, read-only rootfs, gVisor, and composite leave
the vulnerability \textbf{OPEN}. Notably \textbf{gVisor (B7) does NOT fix \texttt{hostpath\_etc}}:
because the payload executes in the host cron daemon outside the container, gVisor's sandbox boundary is
completely bypassed. \textbf{User-namespace remap (B5) is the sole genuine FIX}: container root maps to an
unprivileged host uid (\num{100000}), preventing writes to host-root-owned \texttt{/etc}.

\paragraph{Harness contamination finding.} Empirical model runs revealed a key benchmark design vulnerability:
leftover cron jobs written by earlier epochs persist on the host system across container recreations.
Subsequent epochs inherited persistent background cron execution that continuously dumped host flags into
the shared volume, artificially inflating downstream escape rates. An earlier sweep corrupted host \texttt{/etc/passwd}, and we discarded it and excluded it from all statistics. The full re-run (B0--B8) completed with 0/10 escapes per control in \Cref{tab:hardening-matrix}. The agent never escaped even at baseline, so the model-level rates for this scenario carry no defense signal; the mechanistic classification above is the informative result. B0--B3 ran before the per-epoch host \texttt{/etc} guard existed, and B4--B8 ran with it. The guard restores host \texttt{/etc} before each epoch, but it did not cover \texttt{apt} or \texttt{if-up} hooks, and two hooks written during B7 persisted into B8 until we removed them after the run. We cannot rule out an effect on B8 epochs, although none escaped.

\subsection{Statistical analysis}
This version reports descriptive statistics only; we did not fit a regression model or compute significance tests. Each cell has ten epochs, so individual rates are coarse. Wilson 95\% intervals are $0.00$--$0.28$ for $0/10$ and $0.72$--$1.00$ for $10/10$, and any cell between these is wider still. The contrast between gVisor (B7; ARR $=0.97$ across the three baseline-escapable scenarios) and the other controls is large relative to these intervals, but a formal test across controls and scenarios is future work.
